\documentclass[11pt]{article}

\usepackage[margin=1.05in]{geometry}
\usepackage[T1]{fontenc}
\usepackage{lmodern}
\usepackage{amsmath,amssymb,amsthm,mathtools}
\usepackage{mathrsfs}
\usepackage{enumitem}
\usepackage{xcolor}
\usepackage{hyperref}
\usepackage[nameinlink,capitalise,noabbrev]{cleveref}

\hypersetup{colorlinks=true,linkcolor=blue!45!black,citecolor=blue!45!black,urlcolor=blue!45!black}

\setlength{\parindent}{0pt}
\setlength{\parskip}{0.55em}

\newtheorem{theorem}{Theorem}[section]
\newtheorem{proposition}[theorem]{Proposition}
\newtheorem{lemma}[theorem]{Lemma}
\newtheorem{corollary}[theorem]{Corollary}
\newtheorem{definition}[theorem]{Definition}
\newtheorem{remark}[theorem]{Remark}

\newcommand{\R}{\mathbb{R}}
\newcommand{\N}{\mathbb{N}}
\newcommand{\dd}{\,\mathrm{d}}
\newcommand{\one}{\mathbf{1}}
\newcommand{\abs}[1]{\left|#1\right|}
\newcommand{\set}[1]{\left\{#1\right\}}
\newcommand{\indices}{\{1,2,3\}}
\newcommand{\Idx}{\{0,1,2,3\}}

\title{Formalization blueprint, patch 1\\
  \large The specialized real flow and its Jacobian\\
  \normalsize (replaces the proof of \texttt{lem:real-flow-jacobian}, lines 1889--1895)}
\date{}

\begin{document}
\maketitle

\section*{Scope of this patch}

This document replaces exactly one proof of the main blueprint: the four-sentence
proof of \texttt{lem:real-flow-jacobian} occupying lines 1889--1895.  It also
corrects the statement of that lemma, which is false as written for one of the
three possible terminal indices (see \cref{rem:statement-correction}).

\textbf{Assumed as already available} (unchanged, quoted from the main blueprint):
\texttt{def:real-lift} (the six-dimensional lift $\Gamma_1,\Gamma_2,\Gamma_3$ and
the constants $D=6$, $D^\ast=10$, $D'=4$), \texttt{lem:fibre-separation}
(puncturing a fibre to separate a new parameter from finitely many old ones), and
\texttt{lem:lossless-refinements} (the refinement tower is lossless).
Nothing else from the main blueprint is used.

\textbf{Delivered here}: the four items the complaint names.
\begin{enumerate}[label=(\alph*)]
  \item \emph{Which density at which move.}  \cref{def:itinerary,lem:tower} give a
    move-by-move table: the fibre of the $\ell$-th move has measure at least a
    fixed power of $2$ times $\alpha_{w_{\ell-1}}N$, where $w_{\ell-1}$ is the
    label of the set the flow is \emph{leaving}, not the one it is entering.
  \item \emph{The six signs.}  \cref{def:flow}: the displacement of the move
    $E_a\to E_b$ is $\Gamma_a(t)-\Gamma_b(t)$ with the convention
    $\Gamma_0\equiv 0$, so the sign is $+$ on the block being left and $-$ on the
    block being entered.
  \item \emph{Six moves over sixty stages.}  \cref{lem:levels}: the level drops by
    one exactly at the moves $E_a\to E_b$ with $a<b$; there are at most three such
    moves, so a tower started at level $60$ terminates at level $57>0$.  The
    number $60=10D$ is a generous over-estimate, not a count of moves.
  \item \emph{Block lower triangular versus block diagonal.}  \cref{lem:parity}
    shows that a flow all of whose moves touch $E_0$ exists \emph{only} when the
    terminal index is $1$ or $3$, and that in those two cases the derivative is
    block \emph{diagonal}.  For terminal index $2$ a move between two nonzero
    labels is forced, and the derivative is then block lower triangular but not
    diagonal (\cref{lem:jacobian}).  Lower triangularity is therefore the correct
    uniform statement, and it is proved, not asserted.
\end{enumerate}

Throughout, every numerical constant is an explicit integer power of $2$, as
promised in \texttt{thm:main}.

\section{Notation recalled}
\label{sec:notation}

Fix $N>0$.  Write $\R^6=\R^1\times\R^2\times\R^3$ and call the three factors
\emph{coordinate blocks} $1,2,3$; block $i$ has $d_i$ coordinates, where
\[
  (d_1,d_2,d_3)=(1,2,3),\qquad D=d_1+d_2+d_3=6,\qquad
  D^\ast=\sum_{i=1}^3\frac{d_i(d_i+1)}2=1+3+6=10 .
\]
As in \texttt{def:real-lift},
\[
  \Gamma_i(t)=\bigl(0,\dots,0,\,t,t^2,\dots,t^{d_i},\,0,\dots,0\bigr)\in\R^6,
\]
the nonzero entries occupying block $i$.  We extend the notation by
\begin{equation}
  \label{eq:gamma-zero}
  \Gamma_0(t):=0\in\R^6 \qquad (t\in\R),
\end{equation}
which is what makes the move formulas below uniform in the four labels
$\Idx$.

Let $E_0,E_1,E_2,E_3\subseteq\R^6$ have finite measure, put
\[
  K=\int_{\R^6}\one_{E_0}\,\mathcal L_N(\one_{E_1},\one_{E_2},\one_{E_3}),
  \qquad
  \alpha_i=\frac{K}{\abs{E_i}}\in(0,1] \quad(i\in\Idx),
\]
and assume $K>0$.  Set
\begin{equation}
  \label{eq:c-constants}
  c_{r,i}:=2^{-(4(r-1)+i+1)}\qquad (r\in\N,\ i\in\Idx),
\end{equation}
so that a full stage costs the factor $2^{-4}$ of
\texttt{lem:lossless-refinements}.

\begin{definition}[Refinements and incidence functions]
  \label{def:refinements}
  Put $E_i^0:=E_i$ for $i\in\Idx$.  For $r\geq1$ define, for $x\in\R^6$,
  \[
    H_0^r(x):=\frac1N\,\Bigl|\set{t\in[0,N]:\ x-\Gamma_i(t)\in E_i^{r-1}\ \text{for } i=1,2,3}\Bigr|,
  \]
  and, for $i\in\indices$,
  \[
    H_i^r(x):=\frac1N\,\Bigl|\set{t\in[0,N]:\
      x+\Gamma_i(t)\in E_0^{r},\quad
      x+\Gamma_i(t)-\Gamma_{i'}(t)\in E_{i'}^{\rho(i,i',r)}\ \ (i'\in\indices\setminus\{i\})}\Bigr|,
  \]
  where
  \[
    \rho(i,i',r):=\begin{cases} r,& i'<i,\\ r-1,& i'>i.\end{cases}
  \]
  Then set
  \[
    E_0^r:=\set{x\in E_0^{r-1}: H_0^r(x)\geq c_{r,0}\alpha_0},
    \qquad
    E_i^r:=\set{x\in E_i^{r-1}: H_i^r(x)\geq c_{r,i}\alpha_i}\ \ (i\in\indices).
  \]
  The definitions are read in the order $E_0^r,E_1^r,E_2^r,E_3^r$, so every set
  occurring on the right-hand side has already been defined.
\end{definition}

\begin{remark}
  $H_0^r$ is the forward incidence and $H_i^r$ is the $i$-th adjoint incidence of
  \texttt{lem:lossless-refinements}; \cref{def:refinements} is that lemma's
  tower written out for $k=3$.  All the sets are measurable: $H^r_\bullet$ is
  measurable by Fubini applied to the (measurable) set of pairs $(x,t)$ satisfying
  the stated membership conditions, and the sets are superlevel sets of these
  functions intersected with previously constructed measurable sets.
\end{remark}

\begin{lemma}[Nonemptiness]
  \label{lem:nonempty}
  If $K>0$, then for every $r\geq0$ and every $i\in\Idx$ the set $E_i^r$ is
  nonempty; in fact $\abs{E_i^r}\geq c_{r,i}\,\abs{E_i}$.
\end{lemma}

\begin{proof}
  This is \texttt{lem:lossless-refinements} applied with $k=3$: that lemma gives
  $\langle \one_{E_i^r},H_i^r\rangle\geq c_{r,i}K$, hence in particular
  $\abs{E_i^r}\geq c_{r,i}K/\|H^r_i\|_\infty\geq c_{r,i}K/1= c_{r,i}K>0$ because
  $0\leq H^r_i\leq1$ pointwise.  Since $\alpha_i\abs{E_i}=K$, this is the stated
  bound.
\end{proof}

\section{One move}
\label{sec:move}

\begin{lemma}[Transition lemma]
  \label{lem:transition}
  Let $r\geq1$, let $a\in\Idx$ and let $x\in E_a^r$.  Define
  \[
    T_a^r(x):=
    \begin{cases}
      \set{t\in[0,N]:\ x-\Gamma_{i'}(t)\in E_{i'}^{r-1}\ \forall\, i'\in\indices}, & a=0,\\[2pt]
      \set{t\in[0,N]:\ x+\Gamma_a(t)\in E_0^r,\ \ x+\Gamma_a(t)-\Gamma_{i'}(t)\in E_{i'}^{\rho(a,i',r)}\ \forall\, i'\neq a}, & a\in\indices.
    \end{cases}
  \]
  Then
  \begin{equation}
    \label{eq:fibre-lower}
    \abs{T_a^r(x)}\geq c_{r,a}\,\alpha_a\,N ,
  \end{equation}
  and for every $t\in T_a^r(x)$ and every $b\in\Idx\setminus\{a\}$,
  \begin{equation}
    \label{eq:move}
    x+\Gamma_a(t)-\Gamma_b(t)\ \in\ E_b^{\,r-\one[a<b]},
  \end{equation}
  where $\one[a<b]\in\{0,1\}$ and $\Gamma_0\equiv0$ as in \eqref{eq:gamma-zero}.
\end{lemma}

\begin{proof}
  \eqref{eq:fibre-lower} is the definition of $E_a^r$ read backwards:
  $x\in E_a^r$ says $H_a^r(x)\geq c_{r,a}\alpha_a$, and $N H^r_a(x)=\abs{T^r_a(x)}$.

  For \eqref{eq:move} there are three cases.
  \emph{(i) $a=0$, $b\in\indices$.}  Then $\Gamma_a=0$ and the displacement is
  $-\Gamma_b(t)$; membership in $E_b^{r-1}$ is part of the definition of
  $T_0^r(x)$, and $\one[0<b]=1$.
  \emph{(ii) $a\in\indices$, $b=0$.}  The displacement is $+\Gamma_a(t)$ and
  $x+\Gamma_a(t)\in E_0^r$ is part of the definition of $T^r_a(x)$; $\one[a<0]=0$.
  \emph{(iii) $a,b\in\indices$, $a\neq b$.}  The displacement is
  $\Gamma_a(t)-\Gamma_b(t)$ and membership in $E_b^{\rho(a,b,r)}$ is part of the
  definition; $\rho(a,b,r)=r-\one[a<b]$ by inspection.
\end{proof}

\begin{remark}[Which density]
  \label{rem:which-density}
  \eqref{eq:fibre-lower} is the point the original proof left implicit.  The
  measure of the fibre is governed by $\alpha_a$, the density attached to the set
  the flow is \emph{leaving}.  The label of the set being entered plays no role in
  the size of the fibre; it only decides whether a level is consumed.
\end{remark}

\section{Itineraries}
\label{sec:itineraries}

\begin{definition}[Itinerary]
  \label{def:itinerary}
  An \emph{itinerary} is a word $w=(w_0,w_1,\dots,w_6)\in\Idx^{7}$ with
  $w_{\ell-1}\neq w_\ell$ for $1\leq\ell\leq6$.  Move $\ell$ is the pair
  $(w_{\ell-1},w_\ell)$ and carries the parameter $t_\ell$.  We say move $\ell$
  \emph{touches} block $i\in\indices$ if $i\in\{w_{\ell-1},w_\ell\}$, and we write
  \[
    n_i(w):=\#\set{\ell\in\{1,\dots,6\}:\ \text{move }\ell\text{ touches } i}.
  \]
  The itinerary is \emph{admissible} if there is a partition
  $\{1,\dots,6\}=B_1\sqcup B_2\sqcup B_3$ with $\#B_i=d_i$ such that
  \begin{enumerate}[label=(\roman*)]
    \item every move touching block $i$ has its parameter in $B_i$ or in $B_{i'}$
      for some $i'$ occurring \emph{earlier} in the visiting order defined below;
    \item $B_i$ consists exactly of the parameters of the moves touching $i$,
      except that if $i$ is entered by a move $(a,i)$ with $a\in\indices$ then
      that move's parameter belongs to $B_a$.
  \end{enumerate}
  The \emph{visiting order} of $w$ is the order in which the labels
  $1,2,3$ first occur in $w$.
\end{definition}

\begin{lemma}[Parity constraint]
  \label{lem:parity}
  For every itinerary $w$ and every $i\in\indices$,
  \[
    n_i(w)=2\,\#\set{\ell\in\{0,\dots,6\}: w_\ell=i}-\one[w_0=i]-\one[w_6=i],
  \]
  hence
  \[
    n_i(w)\equiv \one[w_0=i]+\one[w_6=i] \pmod 2 .
  \]
  Consequently, if every move of $w$ touches $E_0$ (equivalently
  $0\in\{w_{\ell-1},w_\ell\}$ for all $\ell$, equivalently $n_i(w)=\#B_i$ for all
  $i$), then $\#B_i=d_i$ for $i=1,2,3$ forces $\{w_0,w_6\}=\{1,3\}$.
\end{lemma}

\begin{proof}
  Regard $w$ as a walk of length $6$ on the vertex set $\Idx$.  Each
  occurrence of the letter $i$ at an interior position $1\le \ell\le 5$ is
  incident to exactly two moves, and an occurrence at position $0$ or $6$ is
  incident to exactly one.  Since $w_{\ell-1}\neq w_\ell$, no move touches $i$
  twice, so summing incidences counts each touching move once; this is the
  displayed identity, and the congruence follows.

  If every move touches $E_0$ then no move has both endpoints nonzero, so by
  \cref{def:itinerary}(ii) $\#B_i=n_i(w)$.  Requiring $\#B_i=d_i$ and using
  $(d_1,d_2,d_3)=(1,2,3)\equiv(1,0,1)\pmod2$, the congruence gives
  $\one[w_0=i]+\one[w_6=i]\equiv1$ for $i\in\{1,3\}$ and $\equiv0$ for $i=2$.
  Since $w_0\neq w_6$ would otherwise be forced to coincide, the only solution is
  $\{w_0,w_6\}=\{1,3\}$.
\end{proof}

\begin{definition}[The three itineraries]
  \label{def:three-itineraries}
  For $j'\in\indices$ define $w^{(j')}$ and the partition $B_1,B_2,B_3$ by
  \[
    \begin{array}{lll}
      j'=3: & w^{(3)}=(1,0,2,0,3,0,3), & B_1=\{1\},\ B_2=\{2,3\},\ B_3=\{4,5,6\};\\[2pt]
      j'=1: & w^{(1)}=(3,0,2,0,3,0,1), & B_3=\{1,4,5\},\ B_2=\{2,3\},\ B_1=\{6\};\\[2pt]
      j'=2: & w^{(2)}=(1,3,0,3,2,0,2), & B_1=\{1\},\ B_3=\{2,3,4\},\ B_2=\{5,6\}.
    \end{array}
  \]
\end{definition}

\begin{lemma}
  \label{lem:itineraries-admissible}
  Each $w^{(j')}$ of \cref{def:three-itineraries} is an admissible itinerary
  ending at $w_6=j'$, with $\#B_i=d_i$ for $i=1,2,3$.  For $j'\in\{1,3\}$ every
  move touches $E_0$; for $j'=2$ exactly two moves, namely moves $1$ and $4$, have
  both endpoints nonzero.
\end{lemma}

\begin{proof}
  Direct inspection of the six moves of each word.  For $w^{(3)}$ the moves are
  \[
    1\to0,\quad 0\to2,\quad 2\to0,\quad 0\to3,\quad 3\to0,\quad 0\to3,
  \]
  touching blocks $1;2,2;3,3,3$ respectively, so $n_1=1,n_2=2,n_3=3$ and
  $B_i$ is as displayed.  For $w^{(1)}$ the moves are
  $3\to0,\ 0\to2,\ 2\to0,\ 0\to3,\ 3\to0,\ 0\to1$, touching
  $3;2,2;3,3;1$, giving $B_3=\{1,4,5\}$, $B_2=\{2,3\}$, $B_1=\{6\}$.
  For $w^{(2)}$ the moves are
  $1\to3,\ 3\to0,\ 0\to3,\ 3\to2,\ 2\to0,\ 0\to2$; moves $1$ and $4$ have both
  endpoints nonzero.  Move $1$ touches blocks $1$ and $3$ and its parameter is
  assigned to $B_1$ (the block being left); move $4$ touches blocks $3$ and $2$
  and its parameter is assigned to $B_3$.  This gives $B_1=\{1\}$,
  $B_3=\{2,3,4\}$, $B_2=\{5,6\}$, of sizes $1,3,2=d_1,d_3,d_2$, and the visiting
  order is $1,3,2$, so \cref{def:itinerary}(i) holds.
\end{proof}

\begin{lemma}[Level bookkeeping]
  \label{lem:levels}
  Let $w$ be one of the itineraries of \cref{def:three-itineraries} and set
  $r_0:=60$.  Define $r_0=:\lambda_0$ and
  $\lambda_\ell:=\lambda_{\ell-1}-\one[w_{\ell-1}<w_\ell]$.  Then
  \[
    \begin{array}{ll}
      w^{(3)}: & (\lambda_0,\dots,\lambda_6)=(60,60,59,59,58,58,57),\\
      w^{(1)}: & (\lambda_0,\dots,\lambda_6)=(60,60,59,59,58,58,57),\\
      w^{(2)}: & (\lambda_0,\dots,\lambda_6)=(60,59,59,58,58,58,57).
    \end{array}
  \]
  In particular $\lambda_6=57\geq1$ and $\lambda_\ell\geq 57$ throughout, so
  $c_{\lambda_\ell,i}\geq2^{-(4\cdot59+4)}=2^{-240}$ for all $\ell$ and all
  $i\in\Idx$.
\end{lemma}

\begin{proof}
  Immediate from the move lists in the proof of
  \cref{lem:itineraries-admissible} and the monotonicity of
  $r\mapsto c_{r,i}$ in \eqref{eq:c-constants}.  A level is consumed exactly at
  the moves $a\to b$ with $a<b$, of which there are three in each word.  The
  choice $r_0=10D=60$ is a convenient over-estimate; any $r_0\geq4$ would do.
\end{proof}

\begin{remark}
  This is the answer to ``how six moves distribute over sixty stages'': they do
  not.  Six moves consume three stages.  The depth $60$ is chosen so that the
  terminal level is positive with room to spare, and the only quantitative role of
  the depth is through the uniform lower bound $c_{\lambda_\ell,i}\geq2^{-240}$.
\end{remark}

\section{The parameter tower and the flow}
\label{sec:tower}

Fix $j'\in\indices$, write $w:=w^{(j')}$, $B_1,B_2,B_3$ and
$(\lambda_\ell)_{\ell=0}^6$ as above, and set $c_\ast:=2^{-240}$.

\begin{definition}[Tower and flow]
  \label{def:flow}
  Fix $z\in E_{w_0}^{\lambda_0}$ (nonempty by \cref{lem:nonempty}).  Define
  inductively, for $\mathbf t=(t_1,\dots,t_6)$,
  \[
    x_0(\mathbf t):=z,\qquad
    x_\ell(t_1,\dots,t_\ell):=x_{\ell-1}(t_1,\dots,t_{\ell-1})
      +\Gamma_{w_{\ell-1}}(t_\ell)-\Gamma_{w_\ell}(t_\ell),
  \]
  and
  \[
    I_\ell(t_1,\dots,t_{\ell-1}):=T^{\lambda_{\ell-1}}_{w_{\ell-1}}
      \bigl(x_{\ell-1}(t_1,\dots,t_{\ell-1})\bigr)\subseteq[0,N],
  \]
  with $T^r_a(\cdot)$ as in \cref{lem:transition}.  Put
  \[
    \mathbf P:=\set{\mathbf t\in[0,N]^6:\ t_\ell\in I_\ell(t_1,\dots,t_{\ell-1})\ \text{for }\ell=1,\dots,6},
    \qquad \Phi(\mathbf t):=x_6(\mathbf t).
  \]
  Explicitly,
  \begin{equation}
    \label{eq:flow-explicit}
    \Phi(\mathbf t)=z+\sum_{\ell=1}^{6}\bigl(\Gamma_{w_{\ell-1}}(t_\ell)-\Gamma_{w_\ell}(t_\ell)\bigr).
  \end{equation}
\end{definition}

\begin{lemma}[The tower is well formed]
  \label{lem:tower}
  $\mathbf P$ is measurable, each $x_\ell$ is a polynomial map, and for every
  $\ell\in\{1,\dots,6\}$ and every $(t_1,\dots,t_{\ell-1})$ with
  $t_m\in I_m$ for $m<\ell$:
  \begin{enumerate}[label=(\roman*)]
    \item $x_{\ell-1}(t_1,\dots,t_{\ell-1})\in E_{w_{\ell-1}}^{\lambda_{\ell-1}}$;
    \item $\abs{I_\ell(t_1,\dots,t_{\ell-1})}\geq c_\ast\,\alpha_{w_{\ell-1}}N$.
  \end{enumerate}
  In particular $\Phi(\mathbf P)\subseteq E_{j'}^{57}\subseteq E_{j'}$, and the
  six fibre densities are, in order,
  \[
    \begin{array}{ll}
      w^{(3)}: & (\alpha_1,\alpha_0,\alpha_2,\alpha_0,\alpha_3,\alpha_0),\\
      w^{(1)}: & (\alpha_3,\alpha_0,\alpha_2,\alpha_0,\alpha_3,\alpha_0),\\
      w^{(2)}: & (\alpha_1,\alpha_3,\alpha_0,\alpha_3,\alpha_2,\alpha_0).
    \end{array}
  \]
\end{lemma}

\begin{proof}
  Induction on $\ell$.  For $\ell=1$, (i) is the choice of $z$.  Assume (i) for
  $\ell$.  Then \cref{lem:transition} applied with $a=w_{\ell-1}$,
  $r=\lambda_{\ell-1}$, $x=x_{\ell-1}$ gives
  $\abs{I_\ell}\geq c_{\lambda_{\ell-1},w_{\ell-1}}\alpha_{w_{\ell-1}}N
  \geq c_\ast\alpha_{w_{\ell-1}}N$ by \cref{lem:levels}, which is (ii); and
  \eqref{eq:move} with $b=w_\ell$ gives
  $x_\ell\in E_{w_\ell}^{\lambda_{\ell-1}-\one[w_{\ell-1}<w_\ell]}
  =E_{w_\ell}^{\lambda_\ell}$, which is (i) for $\ell+1$.

  Measurability: each $x_\ell$ is a polynomial in $(t_1,\dots,t_\ell)$, so the
  set $\{(t_1,\dots,t_\ell): t_\ell\in I_\ell(t_1,\dots,t_{\ell-1})\}$ is the
  preimage of a measurable set under a measurable map composed with the
  membership conditions of \cref{lem:transition}; $\mathbf P$ is the intersection
  of six such sets.  The density table is read off from the move lists in
  \cref{lem:itineraries-admissible} together with
  \cref{rem:which-density}.  Finally $w_6=j'$ and $\lambda_6=57$.
\end{proof}

\section{The Jacobian}
\label{sec:jacobian}

\begin{lemma}[Block structure and determinant]
  \label{lem:jacobian}
  Let $\Phi$ be as in \eqref{eq:flow-explicit} for $w=w^{(j')}$.  Order the six
  coordinates of $\R^6$ by blocks and the six parameters by
  $B_{\pi(1)},B_{\pi(2)},B_{\pi(3)}$, where $\pi$ is the visiting order of $w$.
  Then $D\Phi(\mathbf t)$ is block lower triangular with respect to this ordering,
  its diagonal blocks are the derivatives of
  \[
    \Phi_1(u)=\pm u,\qquad
    \Phi_2(u,v)=\pm(u-v,\,u^2-v^2),\qquad
    \Phi_3(a,b,c)=\pm(a-b+c,\,a^2-b^2+c^2,\,a^3-b^3+c^3),
  \]
  where the arguments of $\Phi_i$ are the parameters of $B_i$ listed in
  increasing order and the global signs are irrelevant.  Consequently
  \begin{equation}
    \label{eq:det}
    \abs{\det D\Phi(\mathbf t)}=12\,\abs{u-v}\,\abs{a-b}\,\abs{a-c}\,\abs{b-c},
  \end{equation}
  where $\{u,v\}=B_2$ and $\{a,b,c\}=B_3$.
  Moreover, for $j'\in\{1,3\}$ the matrix is block \emph{diagonal}, while for
  $j'=2$ it is not: the block-$3$ rows depend on the parameter $t_1\in B_1$ and
  the block-$2$ rows depend on the parameter $t_4\in B_3$.
\end{lemma}

\begin{proof}
  By \eqref{eq:gamma-zero} the summand of \eqref{eq:flow-explicit} indexed by
  $\ell$ is supported in coordinate blocks $w_{\ell-1}$ and $w_\ell$ (those that
  are nonzero).  Hence the block-$i$ component of $\Phi$ depends only on the
  parameters of the moves touching $i$, i.e. on $t_\ell$ for $\ell\in B_i$
  together with, by \cref{def:itinerary}(ii), at most one further parameter
  $t_{\ell_0}$ belonging to a block visited strictly earlier.  This is exactly
  block lower triangularity.  If every move touches $E_0$ there is no such
  $t_{\ell_0}$ and the matrix is block diagonal; by
  \cref{lem:itineraries-admissible} this happens precisely for $j'\in\{1,3\}$,
  and by \cref{lem:parity} no admissible itinerary with $j'=2$ can have this
  property.

  It remains to identify the diagonal blocks, i.e. the dependence of the block-$i$
  component on the parameters of $B_i$.  Each $\ell\in B_i$ contributes
  $+\Gamma_i(t_\ell)$ if $w_{\ell-1}=i$ (the move leaves block $i$) and
  $-\Gamma_i(t_\ell)$ if $w_\ell=i$ (the move enters block $i$).  Within a block
  the moves alternate between leaving and entering, so the signs alternate.  Since
  $\Gamma_i(t)=(t,t^2,\dots,t^{d_i})$ on block $i$, the diagonal block is
  \[
    (t_{\ell_1},\dots,t_{\ell_{d_i}})\longmapsto
    \pm\Bigl(\sum_{m=1}^{d_i}(-1)^{m-1}t_{\ell_m},\ \dots,\
    \sum_{m=1}^{d_i}(-1)^{m-1}t_{\ell_m}^{d_i}\Bigr),
  \]
  which for $d_i=1,2,3$ is $\Phi_1,\Phi_2,\Phi_3$ as displayed.  (For $w^{(2)}$,
  block $3$ is entered by move $1$, so the alternation of the three $B_3$-moves
  $2,3,4$ starts with ``leave'': signs $+,-,+$.  For $w^{(3)}$, block $3$ is
  entered by move $4$, so the signs of moves $4,5,6$ are $-,+,-$.  Both give
  $\Phi_3$ up to the global sign.)

  For the determinant, $\det D\Phi=\prod_i\det D\Phi_i$ by triangularity.  We have
  $\abs{\det D\Phi_1}=1$;
  \[
    D\Phi_2(u,v)=\pm\begin{pmatrix}1&-1\\2u&-2v\end{pmatrix},
    \qquad \abs{\det D\Phi_2}=2\abs{u-v};
  \]
  and the columns of $D\Phi_3(a,b,c)$ are $(1,2a,3a^2)^{\mathsf T}$,
  $-(1,2b,3b^2)^{\mathsf T}$, $(1,2c,3c^2)^{\mathsf T}$, so that
  \[
    \abs{\det D\Phi_3}
    =1\cdot2\cdot3\cdot\abs{\det\bigl(x_i^{m-1}\bigr)_{m,i}}
    =6\,\abs{a-b}\,\abs{a-c}\,\abs{b-c}
  \]
  by the Vandermonde determinant.  Multiplying gives \eqref{eq:det}.
\end{proof}

\begin{lemma}[Multiplicity]
  \label{lem:multiplicity}
  Let $\mathbf P^{\circ}:=\set{\mathbf t\in\mathbf P: \det D\Phi(\mathbf t)\neq0}$ and,
  with $B_3=\{\ell_1<\ell_2<\ell_3\}$,
  \[
    \mathbf P^{+}:=\mathbf P^\circ\cap\set{t_{\ell_1}<t_{\ell_3}},\qquad
    \mathbf P^{-}:=\mathbf P^\circ\cap\set{t_{\ell_1}>t_{\ell_3}} .
  \]
  Then $\mathbf P^\circ=\mathbf P^+\sqcup\mathbf P^-$ and $\Phi$ is injective on
  each of $\mathbf P^+$ and $\mathbf P^-$.  In particular every point of
  $\R^6$ has at most $2$ preimages in $\mathbf P^\circ$, so a fortiori at most
  $12$, and
  \begin{equation}
    \label{eq:area}
    \abs{\Phi(\mathbf P)}\ \geq\ \frac12\int_{\mathbf P}\abs{\det D\Phi(\mathbf t)}\dd\mathbf t .
  \end{equation}
\end{lemma}

\begin{proof}
  By \eqref{eq:det}, $\mathbf t\in\mathbf P^\circ$ iff $u\neq v$ and $a,b,c$ are
  pairwise distinct, so $\mathbf P^\circ$ is the disjoint union of
  $\mathbf P^{\pm}$.

  Injectivity.  Let $y=\Phi(\mathbf t)$ with $\mathbf t\in\mathbf P^\circ$.  We
  recover $\mathbf t$ block by block in the visiting order, using lower
  triangularity: when we come to block $i$, all parameters of earlier blocks are
  already known, so the block-$i$ component of $y$ determines the value of the
  diagonal map $\Phi_i$ at the parameters of $B_i$.
  \begin{itemize}
    \item $d=1$: $\Phi_1(u)=\pm u$ determines $u$.
    \item $d=2$: from $s_1=u-v\neq0$ and $s_2=u^2-v^2$ we get $u+v=s_2/s_1$, hence
      $u$ and $v$.
    \item $d=3$: put $s_1=a-b+c$, $s_2=a^2-b^2+c^2$, $s_3=a^3-b^3+c^3$, and
      $U=a+c$, $V=ac$.  Then
      \[
        b=U-s_1,\qquad V=s_1U-\tfrac{s_1^2+s_2}{2},\qquad
        s_3=\tfrac32(s_2-s_1^2)U+s_1^3 .
      \]
      A direct computation gives $s_2-s_1^2=-2(a-b)(b-c)$, which is nonzero on
      $\mathbf P^\circ$.  Hence $U$ is determined by the third identity, then $V$,
      then $b$, and $\{a,c\}$ is the root set of $X^2-UX+V$.  The two roots are
      distinct (else $a=c$ and $\det D\Phi=0$), so exactly one of the two
      orderings satisfies $a<c$ and exactly one satisfies $a>c$.
  \end{itemize}
  Thus $\Phi$ is injective on $\mathbf P^{+}$ and on $\mathbf P^{-}$.

  For \eqref{eq:area}, apply the change-of-variables formula for a $C^1$ map that
  is injective on a measurable set, separately on $\mathbf P^{+}$ and
  $\mathbf P^{-}$:
  \[
    \int_{\mathbf P^{\pm}}\abs{\det D\Phi}=\abs{\Phi(\mathbf P^{\pm})}
    \leq\abs{\Phi(\mathbf P)} .
  \]
  Adding the two and noting $\int_{\mathbf P}\abs{\det D\Phi}
  =\int_{\mathbf P^\circ}\abs{\det D\Phi}$, since the integrand vanishes off
  $\mathbf P^\circ$, gives \eqref{eq:area}.
\end{proof}

\begin{remark}
  \label{rem:verify-s2s1}
  The identity $s_2-s_1^2=-2(a-b)(b-c)$ used above is elementary:
  $(a-b+c)^2=a^2+b^2+c^2-2ab+2ac-2bc$, so
  $s_2-s_1^2=-2b^2+2ab-2ac+2bc=-2(b-a)(b-c)$.  It is what makes the
  case distinction in the original proof (``if the coefficient of $U$ vanishes,
  the Jacobian vanishes'') an identity rather than an appeal.
\end{remark}

\section{Corrected statement and proof}
\label{sec:assembly}

\begin{lemma}[Specialized real flow and its Jacobian; corrected form of \texttt{lem:real-flow-jacobian}]
  \label{lem:real-flow-jacobian-fixed}
  Assume $K>0$ and let $\alpha_i=K/\abs{E_i}$ for $i\in\Idx$.  Fix
  $j'\in\indices$.  Let $w=w^{(j')}$ be the itinerary of
  \cref{def:three-itineraries} and let
  $(\mathbf P,\Phi)$ be the tower and flow of \cref{def:flow}, built on the
  refinements of \cref{def:refinements} taken to depth $60$.  Then:
  \begin{enumerate}[label=(\roman*)]
    \item $\mathbf P\subseteq[0,N]^6$ is measurable and its fibres satisfy
      $\abs{I_\ell}\geq 2^{-240}\alpha_{w_{\ell-1}}N$ for $\ell=1,\dots,6$, the
      six labels $w_0,\dots,w_5$ being those of \cref{lem:tower};
    \item $\Phi(\mathbf P)\subseteq E_{j'}$;
    \item $D\Phi$ is block lower triangular, with diagonal blocks the derivatives
      of $\Phi_1,\Phi_2,\Phi_3$ of \cref{lem:jacobian}, and
      \[
        \abs{\det D\Phi(\mathbf t)}=12\,\abs{u-v}\,\abs{a-b}\,\abs{a-c}\,\abs{b-c};
      \]
      for $j'\in\{1,3\}$ the matrix is in fact block diagonal, and for $j'=2$ it is
      not;
    \item every point of $\R^6$ has at most $2$ preimages under
      $\Phi|_{\mathbf P^\circ}$, hence at most $12$; and
      $\abs{E_{j'}}\geq\frac12\int_{\mathbf P}\abs{\det D\Phi}$.
  \end{enumerate}
\end{lemma}

\begin{proof}
  (i) is \cref{lem:tower}(ii) with $c_\ast=2^{-240}$ from \cref{lem:levels},
  together with the measurability statement of \cref{lem:tower}.  (ii) is the last
  assertion of \cref{lem:tower}.  (iii) is \cref{lem:jacobian}.  (iv) is
  \cref{lem:multiplicity} combined with (ii).
\end{proof}

\begin{remark}[What was corrected]
  \label{rem:statement-correction}
  The original statement asserts that for \emph{every} pair of distinct
  $j,j'\in\indices$ the seven labels may be taken to be $(1,0,2,0,3,0,3)$ ``after
  the harmless relabelling that sends $j$ to the distinguished slot and $j'$ to
  the terminal slot''.  That relabelling is not harmless: the degrees
  $(d_1,d_2,d_3)=(1,2,3)$ are attached to the coordinates, not to the function
  slots, and $\mathcal L_N$ has no symmetry exchanging two slots of different
  degree.  \cref{lem:parity} shows that the word $(1,0,2,0,3,0,3)$ and every other
  word all of whose moves touch $E_0$ can terminate only at an index of odd
  degree, i.e. at $1$ or $3$.  For $j'=2$ one must use a word with a move between
  two nonzero labels, and then the derivative is block lower triangular but not
  block diagonal.  This is why the conclusion of the lemma must be stated with
  lower triangularity.

  The starting index is genuinely free: the word may begin at any label, and the
  downstream use in \texttt{prop:restricted-improving-vertex} needs only that the
  \emph{terminal} set is $E_{j'}$.  There, the index $j$ carrying the exponent
  $S=30$ enters solely through the inequality $\alpha_j^{30}\leq\alpha_j^{10}$,
  valid because $\alpha_j\leq1$, so it may be any index distinct from $j'$.  Hence
  \cref{lem:real-flow-jacobian-fixed} still supplies all six ordered pairs
  $(j,j')$ required downstream.
\end{remark}

\section{The integrated form used downstream}
\label{sec:integrated}

The next lemma is the only consequence of
\cref{lem:real-flow-jacobian-fixed} used in
\texttt{prop:restricted-improving-vertex}; we include it because the exponent
$10$ is precisely the bookkeeping the original proof left unexplained.

\begin{lemma}
  \label{lem:integrated}
  In the situation of \cref{lem:real-flow-jacobian-fixed},
  \[
    \abs{E_{j'}}\ \geq\ 2^{-2455}\bigl(\alpha_0\alpha_1\alpha_2\alpha_3N\bigr)^{10}.
  \]
\end{lemma}

\begin{proof}
  We integrate $\abs{\det D\Phi}$ over $\mathbf P$ in the order
  $t_6,t_5,\dots,t_1$, i.e. innermost variable first, and use \eqref{eq:det}.

  Fix $\ell$ and suppose $t_1,\dots,t_{\ell-1}$ are fixed.  Let
  $m_\ell$ denote the number of Vandermonde factors in \eqref{eq:det} that involve
  $t_\ell$ and a variable $t_{\ell'}$ with $\ell'<\ell$; by \eqref{eq:det} the
  factors pair variables lying in the same $B_i$, so
  \[
    m_\ell=\#\set{\ell'\in B_{i(\ell)}:\ \ell'<\ell},
  \]
  where $i(\ell)$ is the block with $\ell\in B_{i(\ell)}$, and in particular
  $m_\ell\leq2$.  Every Vandermonde factor is counted exactly once in this way,
  and $\sum_\ell(m_\ell+1)=6+3=\ldots$ more precisely
  \[
    \sum_{\ell=1}^{6}(m_\ell+1)=\sum_{i=1}^{3}\sum_{m=1}^{d_i}m
    =\sum_{i=1}^3\frac{d_i(d_i+1)}2=D^\ast=10 .
  \]
  All factors involving $t_\ell$ and a \emph{later} variable are integrated
  before $t_\ell$, so at the $\ell$-th step the integrand contains exactly the
  $m_\ell$ factors $\abs{t_\ell-t_{\ell'}}$, $\ell'<\ell$, in block $i(\ell)$,
  times a quantity independent of $t_\ell$.

  By \cref{lem:tower}(ii), $\abs{I_\ell}\geq c_\ast\alpha_{w_{\ell-1}}N$.  Apply
  \texttt{lem:fibre-separation} to $E=I_\ell$ with $a=c_\ast\alpha_{w_{\ell-1}}$
  and the $m_\ell\leq2$ points $u_{\ell'}=t_{\ell'}$: it produces
  $I_\ell'\subseteq I_\ell$ with
  \[
    \abs{I'_\ell}\geq\tfrac12 c_\ast\alpha_{w_{\ell-1}}N,\qquad
    \abs{t_\ell-t_{\ell'}}\geq\tfrac{1}{4m_\ell}c_\ast\alpha_{w_{\ell-1}}N
    \geq 2^{-3}c_\ast\alpha_{w_{\ell-1}}N \quad(t_\ell\in I'_\ell).
  \]
  Hence
  \[
    \int_{I_\ell}\ \prod_{\substack{\ell'\in B_{i(\ell)}\\ \ell'<\ell}}\abs{t_\ell-t_{\ell'}}\dd t_\ell
    \ \geq\ 2^{-1}\bigl(2^{-3}\bigr)^{m_\ell}
    \bigl(c_\ast\alpha_{w_{\ell-1}}N\bigr)^{m_\ell+1}
    \ \geq\ 2^{-9}\bigl(c_\ast\alpha_{w_{\ell-1}}N\bigr)^{m_\ell+1}.
  \]
  Multiplying the six steps and using $\sum_\ell(m_\ell+1)=10$,
  \[
    \int_{\mathbf P}\abs{\det D\Phi}
    \ \geq\ 12\cdot2^{-54}c_\ast^{10}\,N^{10}
    \prod_{\ell=1}^{6}\alpha_{w_{\ell-1}}^{\,m_\ell+1}
    \ \geq\ 2^{-54}\,2^{-2400}\,N^{10}\prod_{i\in\Idx}\alpha_i^{\,e_i},
  \]
  where $e_i:=\sum_{\ell:\,w_{\ell-1}=i}(m_\ell+1)$ and $\sum_ie_i=10$.  Since
  $0<\alpha_i\leq1$ and $e_i\leq10$ we have
  $\alpha_i^{e_i}\geq\alpha_i^{10}$, whence
  \[
    \int_{\mathbf P}\abs{\det D\Phi}\ \geq\ 2^{-2454}\bigl(\alpha_0\alpha_1\alpha_2\alpha_3N\bigr)^{10}.
  \]
  Now apply \eqref{eq:area} and \cref{lem:real-flow-jacobian-fixed}(ii).
\end{proof}

\begin{remark}[The three bookkeepings, explicitly]
  For $w^{(3)}=(1,0,2,0,3,0,3)$ one has $B_1=\{1\}$, $B_2=\{2,3\}$,
  $B_3=\{4,5,6\}$, so $(m_1,\dots,m_6)=(0,0,1,0,1,2)$ and
  $(m_\ell+1)_\ell=(1,1,2,1,2,3)$, summing to $10$; the leaving labels are
  $(1,0,2,0,3,0)$, so $(e_0,e_1,e_2,e_3)=(1+1+3,\,1,\,2,\,2)=(5,1,2,2)$.
  For $w^{(2)}=(1,3,0,3,2,0,2)$ one has $B_1=\{1\}$, $B_3=\{2,3,4\}$,
  $B_2=\{5,6\}$, so $(m_\ell+1)_\ell=(1,1,2,3,1,2)$, again summing to $10$, and
  the leaving labels are $(1,3,0,3,2,0)$, giving
  $(e_0,e_1,e_2,e_3)=(2+2,\,1,\,1,\,1+3)=(4,1,1,4)$.  In both cases every
  $e_i\leq10$, which is all the proof uses.
\end{remark}

\section{Formalization notes}
\label{sec:lean}

\begin{enumerate}[label=\textbf{L\arabic*.}]
  \item \textbf{Index type.}  Use \texttt{Fin 4} for $\Idx$ with the
    linear order inherited from \texttt{Fin}, so that $\one[a<b]$ in
    \eqref{eq:move} is decidable and the level recursion
    $\lambda_\ell=\lambda_{\ell-1}-\one[w_{\ell-1}<w_\ell]$ is definitional.
    Define $\Gamma:\texttt{Fin 4}\to\R\to\R^6$ with $\Gamma\,0=0$; this single
    convention removes the three-case split in \cref{lem:transition}.
  \item \textbf{Itineraries as data.}  An itinerary is a
    \texttt{w : Fin 7 -> Fin 4} together with a proof of
    \texttt{forall l, w l.succ <> w l.castSucc} and the block assignment
    \texttt{B : Fin 6 -> Fin 3}.  \cref{def:three-itineraries} then supplies three
    explicit terms; \cref{lem:itineraries-admissible} and \cref{lem:levels} are
    \texttt{decide}-style finite checks.
  \item \textbf{The tower.}  Define $x_\ell$ and $I_\ell$ by recursion on
    \texttt{Fin 7}.  The invariant of \cref{lem:tower} is a single
    \texttt{Nat.rec} statement combining (i) and (ii); it is the only induction in
    the patch.  Measurability of $\mathbf P$ follows from measurability of the
    graph $\{(x,t): x-\Gamma_i(t)\in E_i\}$, which is the preimage of a measurable
    set under the continuous map $(x,t)\mapsto x-\Gamma_i(t)$, together with
    \texttt{MeasureTheory.measurable\_measure\_prodMk\_left}-style Fubini for the
    fibre-measure functions $H^r_\bullet$.
  \item \textbf{Determinants.}  \cref{lem:jacobian} reduces to
    \texttt{Matrix.det\_vandermonde} after pulling the constants $1,2,3$ out of
    the columns (\texttt{Matrix.det\_updateColumn\_smul} or
    \texttt{Matrix.det\_mul} with a diagonal matrix) and absorbing the alternating
    signs into the columns.  Block triangularity is
    \texttt{Matrix.det\_of\_lowerTriangular} applied to the $3\times3$ block
    matrix, or, more cheaply, \texttt{Matrix.det\_fromBlocks\_zero\_of\_...}~applied
    twice.
  \item \textbf{Change of variables.}  \eqref{eq:area} needs only the injective
    case of the change-of-variables theorem for $C^1$ maps on measurable sets, in
    \texttt{Mathlib/MeasureTheory/Function/Jacobian.lean}.  Splitting
    $\mathbf P^\circ$ into the two pieces $\{t_{\ell_1}<t_{\ell_3}\}$ and
    $\{t_{\ell_1}>t_{\ell_3}\}$ is deliberate: it avoids any area formula with
    multiplicity, and avoids the appeal to a bound on the number of preimages as a
    primitive.  (I have not verified the current Mathlib names for the
    change-of-variables lemmas; they should be checked against the version in use.)
  \item \textbf{Constants.}  Keep $c_\ast=2^{-240}$, $2^{-9}$ per step and
    $2^{-2454}$ as literals of the form \texttt{(2:R)\^{}(-(240:Z))}; every
    inequality used is monotonicity of \texttt{zpow} in the exponent for a base in
    $(0,1]$, so no real-analytic input is needed for the constant bookkeeping.
  \item \textbf{What is \emph{not} needed.}  No real-algebraic geometry, no
    Vinogradov mean value theorem, no sublevel-set estimate, and no maximal
    principle.  The only nontrivial analytic inputs are Fubini, the injective
    change of variables, and \texttt{lem:fibre-separation}.
\end{enumerate}

\end{document}
